\subsection{2 预训练}\label{pretraining}

MAI-Thinking-1 的基础模型 MAI-Base-1 是一个 35B-active / 1T-total 的稀疏 MoE 模型，从零开始在 Azure 平台内由 Microsoft 运营的集群上的 8K GB200 GPU 上预训练，使用自研分布式训练基础设施（第 2.8 节）。我们的训练语料完全由内部构建，来源包括公开可获取的数据与获取的数据，不含任何蒸馏数据（第 2.4 节）。基于该语料，我们进行了 30 万亿 token 的主预训练阶段，随后是总计 3.55 万亿 token 的中期训练阶段（第 2.6 节）。

我们在预训练中的目标是产出一个在广泛领域内都具有能力的基础模型，且其能力源于人类知识，而非对现有 AI 模型的模仿。我们建立了一个消融协议，通过沿着一条由规模递增的预训练模型构成的“阶梯”测量效率增益（第 2.2 节），仔细考察我们模型决策的扩展性质。为支持这些消融实验，我们开发了一套预训练指标（第 2.3 节），包括覆盖广泛学科、任务和数据分布的留出负对数似然（NLL）评估。我们报告 MAI-Base-1 与同期基础模型的性能，以说明这些评估在未经后训练（post-training）的模型上的应用（第 2.7 节）。

\pandocbounded{\includegraphics[keepaspectratio,alt={image}]{images/f18ba1d671e0ba94891deedd962868736db005a9a9906241919b4a31e6758065.jpg}}

图 2. MAI-Base-1 架构总览。左：Transformer 主体的整体布局，我们将高稀疏度 MoE 层与小型稠密 FFN 交错，并将全局注意力与局部注意力交错。右：MoE 层，在压缩潜在空间中每个 token 激活 512 个专家中的 8 个。



\subsection{2.1 模型架构}\label{model-architecture}

MAI-Base-1 的模型架构在设计上充分利用了规模在我们的基础设施和硬件上的优势。我们使用仅解码器（decoder-only）Transformer（Vaswani et al., 2017），包含周期性交替的局部与全局注意力层，以及稠密与 MoE 前馈块。架构的高层总览如图 2 所示。在每一层中，我们在残差相加之前的输入和输出处都使用 RMSNorm（Zhang and Sennrich, 2019）（Gemma Team et al., 2025）。模型中不使用任何偏置（bias），并且输入与输出 embedding 权重共享（tied）。虽然我们观察到自研训练的 tokenizer 能带来一些模型性能提升，但我们选择了词汇量为 200,019 的 o200k\_base tokenizer（OpenAI, 2022; Hurst et al., 2024），以简化与现有自研工具和工作流的集成。我们架构设计中的一个关键考量是与底层 GPU 基础设施的协同优化，以确保预训练和强化学习（RL）都能高效扩展。

注意力机制。我们遵循 Gemma 3（Gemma Team et al., 2025）的周期性注意力设计，将 5 个局部注意力层（Beltagy et al., 2020）与 1 个全局注意力层配对，这显著降低了训练的注意力计算成本以及推理期间的 KV 缓存大小。局部注意力层使用旋转位置编码（RoPE）（Su et al., 2023），滑动窗口大小为 512，基准频率为 10,000。全局注意力层不使用任何位置编码（Kazemnejad et al., 2023），其性能与 RoPE 相当，同时更高效。我们使用分组查询注意力（GQA），含 8 个 KV 头（Ainslie et al., 2023），每头维度为 128。我们对查询和键应用 RMSNorm（Roy et al., 2020; Henry et al., 2020; Dehghani et al., 2023; Wortsman et al., 2023）。这些标准组件的使用使我们能够在长上下文上利用 FlashAttention-4（Zadouri et al., 2026）和 Ulysses 风格上下文并行（Jacobs et al., 2023）实现高效训练。

前馈网络与 MoE。对于每个块中的前馈层，我们在 MoE 层与稠密前馈网络（FFN）之间交替。我们发现，将高稀疏度层与零稀疏度（稠密）层配对的这种模式，其扩展性与均衡稀疏度分配---即全程使用中稀疏度 MoE---相当，而在 iso-active 与 iso-total 参数量设置下墙钟时间更高效（第 2.2.3 节）。我们还发现，全层 MoE 布局更依赖共享专家，而向交错布局添加共享专家几乎无益。我们的第一个前馈层是稠密的（DeepSeek-AI et al., 2025a; Kimi Team et al., 2026; Dai et al., 2024），并且我们对稠密和 MoE 前馈层都使用 SwiGLU（Shazeer, 2020）。

我们采用 LatentMoE 设计（NVIDIA et al., 2025），在 all-to-all 调度之前应用共享降维投影。在 all-to-all 聚合之后，潜在表示被投影回其原始维度。重要的是，路由决策基于原始表示，每个压缩表示通过 softmax 门控路由到 512 个专家中的 8 个。

我们使用全局批次负载均衡损失（Qiu et al., 2025; Yang et al., 2025），其中经验专家频率通过对同一全局批次内的数据并行 worker 和微批次进行聚合获得。我们发现聚合策略比负载均衡损失类型重要得多。只要保证全局聚合，我们使用的 GShard 式损失（Lepikhin et al., 2020）表现与无辅助损失变体（Wang et al., 2024a）相似。

值得注意的是，我们观察到关于最优负载均衡方法的实验结果可能取决于专家容量。在有限容量下，即使精心设计 token 丢弃策略以避免细微的因果泄漏，并且在充裕的容量预算下保持较低的丢弃率，所得结论仍可能与无丢弃设置下得出的结论不同。因此我们最终采用完全无丢弃（dropless）的 MoE 实现（详见第 2.8.1 节），它支持 all-to-all 通信的可变消息大小，并保证在高失衡情况下内存使用有界。



表 1. MAI-Base-1 架构家族中各模型规模的配置。(Expert-)FFN 大小指第一个线性层的输出大小，由于 SwiGLU，它是第二个线性层输入大小的 2 倍。MAI-Base-1 的配置与 L78 略有不同，以提升训练与推理效率。

{\def\LTcaptype{none} % do not increment counter
\begin{longtable}[]{@{}llllllllll@{}}
\toprule\noalign{}
\endhead
\bottomrule\noalign{}
\endlastfoot
\multirow{2}{*}{模型} & \multirow{2}{*}{激活参数} &
\multirow{2}{*}{总参数} & \multirow{2}{*}{层数} &
\multirow{2}{*}{隐藏维度} & \multirow{2}{*}{FFN} & \multirow{2}{*}{降维投影} & \multirow{2}{*}{专家 FFN} & \multirow{2}{*}{Top-k/专家数} &
\multirow{2}{*}{KV/Q} \\
 \\
L12 & 365M & 3.9B & 12 & 1024 & 2048 & 512 & 1536 & 8/512 & 8/16 \\
L18 & 760M & 13B & 18 & 1536 & 3072 & 768 & 2304 & 8/512 & 8/16 \\
L24 & 1.5B & 30B & 24 & 2048 & 4096 & 1024 & 3072 & 8/512 & 8/24 \\
L30 & 2.6B & 58B & 30 & 2560 & 5120 & 1280 & 3840 & 8/512 & 8/32 \\
L36 & 4.0B & 100B & 36 & 3072 & 6144 & 1536 & 4608 & 8/512 & 8/32 \\
L42 & 6.1B & 159B & 42 & 3584 & 7168 & 1792 & 5376 & 8/512 & 8/40 \\
L66 & 21.7B & 615B & 66 & 5632 & 11264 & 2816 & 8448 & 8/512 & 8/64 \\
L78 & 35.6B & 1015B & 78 & 6656 & 13312 & 3328 & 9984 & 8/512 & 8/80 \\
MAI-Base-1 & 34.7B & 962B & 78 & 6656 & 13312 & 3072 & 10240 & 8/512 &
8/80 \\
\end{longtable}
}

\subsection{2.2 模型消融方法论}\label{model-ablation-methodology}

下面我们描述用于评估预训练中不同架构与数据相关设计选择的实验设置。一个关键的考量是规模的作用：随着计算预算增加，模型创新或新数据集带来的表面收益往往会减弱。因此，在将改动纳入我们的基线之前，我们必须在不同规模下收集充分的支撑证据。在接下来的小节中，我们介绍使我们能够以有原则且可复现的方式执行这一过程的主要组件：规模阶梯与效率增益。

\subsection{2.2.1 规模阶梯}\label{scaling-ladder}

所有消融实验都明确以扩展性质为核心关注点。为此，我们利用规模阶梯：对于任何消融实验，我们以恒定的每（激活）参数 token 数（TPP）训练不同规模的模型，并将扩展曲线与我们的基线进行比较。TPP 的选择取决于实验的性质。例如，大多数架构消融实验在 Chinchilla 最优区域（Hoffmann et al., 2022）附近以 100-200 TPP 进行，而我们的主运行通常在更高的 500-1,000 TPP 下过度训练，以得到相对更紧凑、适合重度推理需求的模型。消融实验使用最新的阶梯进行，被接受的改动成为下一版基线阶梯配置的一部分，如此往复。

对于第 2.1 节所述的架构家族，一个模型配置由单一参数完全确定：层数 L。对于我们所使用的 5:1 局部与全局注意力之比，L 必须是 6 的倍数，以便最后一个注意力层应用全局注意力。隐藏维度 D 则由 D =
\(L \times \frac { 2 5 6 } { 3 }\) 确定，以保持阶梯上相同的纵横比。这一纵横比选择与一些近期模型（OpenAI, 2025; Liu et al., 2024; Gemma Team, 2024）类似，并且由于 L 始终是 6 的倍数，在我们的设置中会得到对硬件更友好的 D 值。查询头数量设为 L，并向上取整到最近的 16 的倍数，以便于使用张量并行进行推理服务。每个稠密前馈层将隐藏维度扩展 2 倍。对于 LatentMoE 层，我们使用 2$\times$ 的压缩因子，并在每个专家内部将压缩表示扩展 3$\times$ 倍。

完整规模阶梯的细节见表 1。MAI-Base-1 的最终架构有 78 层。虽然它基于阶梯中的 L78，我们略微调整了其专家输入大小以适配 DeepEP（Zhao et al., 2025a），DeepEP 要求 all-to-all 集合通信的隐藏维度可被 512 整除。



表 2. 与 MAI-Base-1 的交错布局（8/512 MoE + 稠密 FFN）相比，两种全层 MoE 变体的效率增益。以
\(\mathrm { E G } _ { \mathrm { F L O P s } }\) 衡量时两种变体都接近基线，但在
\(\mathrm { E G } _ { \mathrm { T i m e } }\) 上明显落后。

{\def\LTcaptype{none} % do not increment counter
\begin{longtable}[]{@{}lllll@{}}
\toprule\noalign{}
\endhead
\bottomrule\noalign{}
\endlastfoot
\multirow{2}{*}{评估类别} & \multicolumn{2}{l}{%
全层 MoE（8/384）} & \multicolumn{2}{l@{}}{%
全层 MoE（7+1 共享/384）} \\
& \textbackslash mathtt \{ E G \} \_ \{ \textbackslash mathtt \{ F L O P
s \} \} & \textbackslash mathbf \{ E G \} \_ \{ \textbackslash mathrm \{
T i m e \} \} & \textbackslash tt E G \_ \{ \textbackslash tt F L O P s
\} & \textbackslash mathbf \{ E G \} \_ \{ \textbackslash mathrm \{ T i
m e \} \} \\
General & 0.88 & 0.69 & 0.99 & 0.78 \\
STEM & 0.93 & 0.73 & 1.04 & 0.82 \\
Math & 0.94 & 0.74 & 1.06 & 0.84 \\
Code & 0.94 & 0.74 & 1.03 & 0.82 \\
Multilingual & 0.94 & 0.75 & 1.05 & 0.82 \\
Weighted Average & 0.94 & 0.73 & 1.03 & 0.82 \\
\end{longtable}
}

\subsection{2.2.2 效率增益}\label{efficiency-gain}

在我们的架构与扩展讨论中，效率增益（EG）衡量的是：基线模型需要多高的训练成本才能达到给定候选模型达到的评估指标（Kaplan et al., 2020; Hoffmann et al., 2022）。大多数情况下，评估指标是基于我们的预训练评估套件（第 2.3 节）定义的聚合评估损失。我们首先将缩放定律拟合到基线阶梯上，

\[
L = f ( C ) = A C ^ { - \alpha } + E\tag{1}
\]

其中 \(C\) 表示训练成本，如给定基础设施下的 FLOPs 或时间。A 是控制可约损失项整体幅度的缩放系数，$\alpha$ 是控制损失随计算量增加而下降速度的缩放指数，E 是不可约损失。对于以成本 \(C ^ { \prime } { _ { ; } }\) 达到 \(L ^ { \prime }\) 的候选运行，我们计算达到相同损失所需的基线成本预算
\(f ^ { - 1 } ( L ^ { \prime } )\)，并将 EG 定义为

\[
\operatorname { E G } = { \frac { f ^ { - 1 } ( L ^ { \prime } ) } { C ^ { \prime } } }\tag{2}
\]

例如，$\mathrm { E G } \textasciitilde{} = \textasciitilde{} 1
. 3 $ 意味着基线模型需要高出 30\% 的成本才能达到候选模型的损失（Liu et al., 2025a; Dey et al., 2023）。我们还在不同的成本预算下进行消融实验，以获得一系列 EG vs.~cost（EG 随成本变化）曲线，这有助于我们找到随规模增加而持续存在的改进。

我们最常用 FLOPs 来定义成本 \(C ,\)，这有意不考虑训练中可实现的墙钟效率。成熟架构的实现往往随时间积累了大量的性能优化，这可能使它们看起来比更新的变体更强。通过将指标与 FLOPs 利用率（MFU）（Chowdhery et al., 2023）解耦，我们可以考察在对其实现投入同等优化努力的情况下，不同架构之间的相对表现。然而，在考虑到我们现有技术栈确实关心硬件效率的情况下，我们使用上述时间定义作为成本 C。在这些情况下，我们明确将效率增益称为 \(\mathrm { E G } _ { \mathrm { T i m e } }\)。

\subsection{2.2.3 模型消融示例}\label{model-ablation-examples}

我们在日常消融实验中广泛使用阶梯与 EG 比较。下面我们展示两组实验，说明 EG 如何在不同评估类别间变化，以及 EG 如何扩展以反映与我们基础设施的协同设计因素。

示例 1：稀疏度分配。表 2 将我们的交错布局（高稀疏度 MoE + 稠密 FFN）与更常见的全程使用中稀疏度 MoE 层的设计进行比较。实验在稀疏度受控的设置下进行，其中所有模型家族共享相似的总参数与激活参数之比；这是通过固定 top-k（8，或 7+1 带一个共享专家（Dai et al., 2024; Qwen et al., 2025）），同时调整全层 MoE 家族的专家扩展率和专家数量实现的。我们报告沿两个成本维度定义的 EG：FLOPs 和训练时间，在第 2.2.1 节所述的阶梯模型 L12--30 上测量。虽然全层 MoE 结合共享专家在 FLOPs 上达到了 1.03 的 EG，但一旦考虑训练 MFU，交错范式提供了更好的权衡，这反映在
\(\mathrm { E G } _ { \mathrm { T i m e } } < 1\)



Top 8 / 256 专家 Top 8 / 512 专家 Top 8 / 1024 专家

\pandocbounded{\includegraphics[keepaspectratio,alt={image}]{images/1daf88df3f0054dd77d341abff6f04e889eda8a2292d82b17f5890ee45a0b50b.jpg}}

图 3. 随模型家族稀疏度增加，各评估任务类别上的预训练效率增益（第 2.2.2 节）。结果为针对每个稀疏度水平训练一条模型阶梯（第 2.2.1 节）获得。由于每个数据点都与缩放定律（式 1）进行比较，基线原始数据点的 EG 可能不严格等于 1。

示例 2：稀疏度扩展。可扩展的架构设计应在稀疏度增加时带来健康的改进。图 3 表明我们的架构家族具备这一特性，随着专家数量从 256 增加到 1024（top-k = 8，MoE 稀疏度水平从 32$\times$ 到 128$\times$），效率增益保持一致。对于 MAI-Base-1，我们选择 top-8/512 配置以平衡模型质量与训练和推理效率，并将稀疏度扩展的进一步探索留给未来工作。
